\section{Definición del contrato asiático}
\label{sec:contrato_asiatico}

El instrumento estudiado es una call asiática europea de precio medio, cuyo
payoff depende de la media aritmética de $N$ precios observados en un calendario
de fixings $0<t_1<\cdots<t_N=T$ \parencite{kemna_vorst_1990,turnbull_wakeman_1991}:
\begin{equation}
    \bar P_T=\frac{1}{N}\sum_{i=1}^{N}P_{t_i},
    \qquad
    H_T=(\bar P_T-K)^+ .
    \label{eq:asian_payoff}
\end{equation}

En una fecha intermedia $t<T$, sea $m(t)$ el número de fixings ya observados. Se
define su contribución al promedio terminal como
\begin{equation}
    C_t=\frac{1}{N}\sum_{i=1}^{m(t)}P_{t_i},
    \qquad
    \bar P_T=C_t+\frac{1}{N}\sum_{i=m(t)+1}^{N}P_{t_i}.
    \label{eq:asian_average_decomposition}
\end{equation}
Así, el spot actual no basta para valorar el contrato: el estado debe recoger
también la parte del promedio ya fijada. En los experimentos se utiliza un
vencimiento máximo de un año, $N=12$ fixings y strike normalizado $K=1$. Se
denota $\tau=T-t$ y las variables monetarias se expresan respecto a $K$.

Esta estructura es coherente con la motivación económica del trabajo: si un
usuario adquiere capacidad de cómputo de forma recurrente, su exposición puede
estar más relacionada con el coste medio soportado durante el periodo que con un
único precio terminal. Todos los experimentos utilizan nocional unitario; los
resultados pueden escalarse linealmente a otras cantidades contractuales bajo los
supuestos del modelo.


\section{Problema de valoración y medida de valoración}
\label{sec:medida_q}

Para un tipo de interés determinista $r$, el valor modelizado se define como
\begin{equation}
    V_t^{\mathbb Q}
    =e^{-r\tau}\,
    \mathbb E^{\mathbb Q}
    \left[(\bar P_T-K)^+\mid\mathcal F_t\right].
    \label{eq:asian_pricing_q}
\end{equation}
La notación $\mathbb Q$ distingue la dinámica empleada para valorar de la
dinámica física $\mathbb P$, siguiendo la separación habitual entre
probabilidades históricas y de valoración \parencite{uned_mod7_2026}. Sin
embargo, la interpretación requiere cautela. La capacidad de cómputo es un
servicio no almacenable y el precio spot utilizado no corresponde a un activo
que pueda comprarse y almacenarse mediante una estrategia estándar de
\textit{cash-and-carry}. Por ello, no se identifica empíricamente una medida
martingala equivalente única a partir de un mercado completo de derivados. Esta
limitación es análoga a la que aparece en otros mercados de bienes no
almacenables \parencite{bessembinder_lemmon_2002,bandi_su_2026}.

Bajo $\mathbb P$ se estima una dinámica Log-OU a partir de la serie H100,
\begin{equation}
    dX_t=\kappa_P(\theta_P-X_t)dt+\sigma_PdW_t^{\mathbb P},
    \label{eq:logou_p}
\end{equation}
mientras que para valoración se utiliza una ley de la misma familia,
\begin{equation}
    dX_t=\kappa_Q(\theta_Q-X_t)dt+\sigma_QdW_t^{\mathbb Q}.
    \label{eq:logou_q}
\end{equation}
Los parámetros $(\kappa_Q,\theta_Q,\sigma_Q)$ deben interpretarse, por tanto,
como parámetros de una \emph{ley de valoración modelizada}. El escenario central
$Q_0$ adopta la calibración física como referencia,
\begin{equation}
    Q_0:\quad
    \kappa_Q=\widehat\kappa_P,\qquad
    \theta_Q=\widehat\theta_P,\qquad
    \sigma_Q=\widehat\sigma_P,
    \label{eq:q0}
\end{equation}
igualando por construcción las leyes física y de valoración del índice en el
escenario central, sin introducir una prima de riesgo en esa parametrización.
Esta elección no constituye evidencia de ausencia de prima de riesgo en el mercado real. El experimento amplía posteriormente estos parámetros dentro de bandas
de diseño para evitar que los surrogates aprendan exclusivamente una calibración
puntual. En consecuencia, los precios obtenidos son precios de modelo
condicionados a la especificación de $\mathbb Q$, no cotizaciones observadas de
opciones H100.


\section{Geometric Brownian Motion como benchmark}
\label{sec:gbm_benchmark}

El movimiento browniano geométrico (GBM) sirve como benchmark clásico
\parencite{black_scholes_1973,merton_1973,uned_mod7_2026}. Su dinámica convencional
de valoración es
\begin{equation}
    dP_t=rP_tdt+\sigma P_tdW_t^{\mathbb Q},
    \label{eq:gbm_q}
\end{equation}
cuya transición exacta es
\begin{equation}
    P_{t+u}=P_t\exp\left[\left(r-\frac{1}{2}\sigma^2\right)u
    +\sigma\sqrt{u}\,Z\right],\qquad Z\sim\mathcal N(0,1).
    \label{eq:gbm_exact}
\end{equation}
La opción aritmética se valora por simulación, pues su payoff depende de varios
fixings y no dispone en general de una solución cerrada tan simple como la del
caso europeo estándar \parencite{kemna_vorst_1990,glasserman_2004}. El GBM se
utiliza únicamente como benchmark clásico sin reversión a la media; la evidencia
H100 y el objetivo económico del TFM motivan que el Log-OU sea la especificación
principal.


\section{Modelo Log-OU y reversión a la media}
\label{sec:logou}

El modelo principal sigue una estructura de Ornstein--Uhlenbeck para el
log-precio, inspirada en el modelo de un factor de \textcite{schwartz_1997}. Sea
$Y_t=\log P_t$. Bajo la medida física,
\begin{equation}
    dY_t=\kappa_P(\vartheta_P-Y_t)dt+\sigma_PdW_t^{\mathbb P}.
    \label{eq:logou_physical_unscaled}
\end{equation}
Para trabajar en unidades relativas al strike se define
$X_t=\log(P_t/K)$ y $\theta_P=\vartheta_P-\log K$. La normalización no altera
$\kappa_P$ ni $\sigma_P$.

Para parámetros genéricos $(\kappa,\theta,\sigma)$, la transición exacta en un
intervalo $\Delta$ es
\begin{equation}
    X_{t+\Delta}=\theta+e^{-\kappa\Delta}(X_t-\theta)
    +\sigma\sqrt{\frac{1-e^{-2\kappa\Delta}}{2\kappa}}\,Z,
    \qquad Z\sim\mathcal N(0,1).
    \label{eq:logou_exact_transition}
\end{equation}
Por tanto,
\begin{equation}
    \mathbb E[X_{t+u}\mid X_t]
    =\theta+e^{-\kappa u}(X_t-\theta),
    \qquad
    \operatorname{Var}(X_{t+u}\mid X_t)
    =\frac{\sigma^2}{2\kappa}(1-e^{-2\kappa u}).
    \label{eq:logou_conditional_mean}
\end{equation}
La desviación respecto al nivel de largo plazo decae exponencialmente y la vida
media es $t_{1/2}=\log(2)/\kappa$. Para $\kappa>0$, la distribución estacionaria
del log-precio es normal con varianza $\sigma^2/(2\kappa)$. El modelo garantiza
precios positivos, permite simulación exacta y mantiene una parametrización
parsimoniosa, ventajas importantes dada la longitud limitada de la serie H100.
No obstante, impone parámetros constantes, gaussianidad y un único factor, por
lo que no incorpora saltos, volatilidad estocástica, estacionalidad ni cambios
de régimen.

La simulación bajo $\mathbb Q$ utiliza la misma transición exacta sustituyendo
los parámetros físicos por $(\kappa_Q,\theta_Q,\sigma_Q)$ y recupera el precio
como $P_{t+\Delta}=K\exp(X_{t+\Delta})$. Esta formulación permite estudiar una
familia controlada de escenarios de valoración sin presentar sus parámetros
como una calibración neutral al riesgo observada en mercado.


\section{Calibración física y construcción del dominio paramétrico}
\label{sec:logou_calibration}

La dinámica física se estima mediante la representación AR(1) implícita en la
transición exacta del OU. Para observaciones equiespaciadas,
\begin{equation}
    Y_t=a+\phi Y_{t-\Delta}+\varepsilon_t,
    \qquad
    \phi=e^{-\kappa_P\Delta},
    \qquad
    a=\vartheta_P(1-\phi).
    \label{eq:ar1_logou}
\end{equation}
De aquí se obtienen
\begin{equation}
    \widehat\kappa_P=-\frac{\log\widehat\phi}{\Delta},\qquad
    \widehat\vartheta_P=\frac{\widehat a}{1-\widehat\phi},\qquad
    \widehat\sigma_P=\widehat\sigma_\varepsilon
    \sqrt{\frac{2\widehat\kappa_P}{1-\widehat\phi^2}}.
    \label{eq:ou_ar1_mapping}
\end{equation}
Se adopta $\Delta=1/365$, por lo que $\kappa_P$ y $\sigma_P$ quedan expresados
en escala anual.

Dada la corta historia disponible, la incertidumbre de calibración se analiza
mediante un bootstrap residual del AR(1) con $20\,000$ réplicas
\parencite{efron_tibshirani_1993}. En cada réplica se remuestrean los residuos
centrados, se reconstruye recursivamente una serie y se reestiman los parámetros;
se conservan únicamente realizaciones con $0<\widehat\phi^*<1$. El bootstrap se
utiliza para estudiar la sensibilidad de la calibración y para definir el dominio
del surrogate, no como una identificación de parámetros neutrales al riesgo.

Las bandas principales de diseño para la velocidad de reversión y la volatilidad
corresponden aproximadamente a los percentiles P10--P90 del bootstrap,
\begin{equation}
    \kappa_Q\in[25,87],\qquad \sigma_Q\in[0{,}55,0{,}71].
    \label{eq:kappa_design_band}
\end{equation}
El nivel de largo plazo se trata separadamente para no confundir incertidumbre
física con ajuste bajo $\mathbb Q$. Con
$K_{\mathrm{ref}}=\exp(\widehat\vartheta_P)$ y el estado normalizado respecto a
ese nivel, se fija el rango
\begin{equation}
    \theta_Q\in[-0{,}1641,0{,}1610],
    \label{eq:theta_design_band}
\end{equation}
con $\theta_Q=0$ en el escenario central. Estas bandas son dominio de diseño,
no intervalos de confianza de una medida neutral al riesgo.

Como diagnóstico adicional, la capacidad predictiva a un paso del AR(1) se
compara fuera de muestra con un paseo aleatorio. Esta comparación no pretende
probar que el Log-OU sea la ley verdadera del mercado, sino comprobar si la
reversión a la media aporta información frente a un benchmark persistente.


\section{Randomized Quasi-Monte Carlo}
\label{sec:rqmc}

La valoración requiere integrar sobre los fixings futuros de una opción
path-dependent. Para un estado $\mathbf x$ y $d$ fixings pendientes,
\begin{equation}
    V(\mathbf x)=\mathbb E^{\mathbb Q}[g(\mathbf Z;\mathbf x)]
    =\int_{[0,1]^d}g(\Phi^{-1}(\mathbf u);\mathbf x)\,d\mathbf u.
    \label{eq:unit_cube_integral}
\end{equation}
El Monte Carlo convencional se sustituye por redes Sobol aleatorizadas mediante
\textit{scrambling}, es decir, Randomized Quasi-Monte Carlo (RQMC)
\parencite{sobol_1967,owen_1997}. La aleatorización conserva la cobertura de
baja discrepancia y permite disponer de réplicas independientes para medir la
variabilidad del estimador. No se impone una tasa asintótica específica para el
payoff asiático; la elección de RQMC se apoya en su comportamiento numérico y en
su capacidad para generar etiquetas precisas a coste controlado.

Para cada punto Sobol se transforman sus coordenadas a normales estándar y se
propagan secuencialmente mediante la transición exacta Log-OU. Si los $d$
fixings futuros son $t_{j_1},\ldots,t_{j_d}$,
\begin{equation}
    \bar P_T^{(n)}=C_t+\frac{1}{N}\sum_{\ell=1}^{d}P_{t_{j_\ell}}^{(n)},
    \qquad
    \widehat V_{RQMC}=e^{-r\tau}\frac{1}{M}
    \sum_{n=1}^{M}(\bar P_T^{(n)}-K)^+.
    \label{eq:rqmc_asian_price}
\end{equation}
La transición exacta evita error de discretización Euler; el error numérico
procede esencialmente de la integración.

Cada estado utiliza una aleatorización propia. Los conjuntos de entrenamiento,
validación y test emplean $M=8192=2^{13}$ puntos y una aleatorización por estado.
El conjunto de referencia emplea 16 aleatorizaciones independientes de 8192
puntos, de modo que cada etiqueta de referencia agrega $131\,072$ evaluaciones
de trayectoria. Conviene distinguir este RQMC de un segundo uso de Sobol: el
diseño de escenarios sintéticos descrito más adelante. El primero integra una
esperanza condicional; el segundo cubre el espacio de estados.


\section{Estimación pathwise de Delta}
\label{sec:pathwise_delta}

La sensibilidad principal es
\begin{equation}
    \Delta_t=\frac{\partial V_t}{\partial P_t},
    \label{eq:delta_definition}
\end{equation}
manteniendo fija la contribución histórica $C_t$. Bajo el Log-OU,
condicionada una realización de las innovaciones,
\begin{equation}
    \frac{\partial P_{t+u}}{\partial P_t}
    =e^{-\kappa_Q u}\frac{P_{t+u}}{P_t}.
    \label{eq:logou_spot_pathwise}
\end{equation}
Como los fixings pasados no cambian ante una perturbación del spot actual,
$\partial C_t/\partial P_t=0$ y
\begin{equation}
    \frac{\partial\bar P_T}{\partial P_t}
    =\frac{1}{N}\sum_{i=m(t)+1}^{N}
    e^{-\kappa_Q(t_i-t)}\frac{P_{t_i}}{P_t}.
    \label{eq:asian_average_pathwise_derivative}
\end{equation}
Por tanto, casi seguramente,
\begin{equation}
    \Delta_t=e^{-r\tau}\mathbb E^{\mathbb Q}\left[
    \mathbf 1_{\{\bar P_T>K\}}
    \frac{1}{N}\sum_{i=m(t)+1}^{N}
    e^{-\kappa_Q(t_i-t)}\frac{P_{t_i}}{P_t}
    \;\middle|\;\mathcal F_t\right].
    \label{eq:asian_logou_delta}
\end{equation}
La misma trayectoria RQMC proporciona así el precio y su derivada, sin
revaloraciones adicionales. Esta propiedad es especialmente útil para DML.
Además, al normalizar $s=P_t/K$ y $v=V_t/K$ se mantiene
$\partial v/\partial s=\Delta_t$.

La implementación pathwise se valida mediante diferencias finitas centrales con
\textit{common random numbers},
\begin{equation}
    \Delta_t^{FD}(h)=\frac{V(P_t+h,C_t,\ldots)-V(P_t-h,C_t,\ldots)}{2h},
    \label{eq:central_fd_delta}
\end{equation}
utilizadas exclusivamente como diagnóstico. La Delta debe interpretarse como
sensibilidad del modelo; no implica por sí misma que el spot de cómputo sea un
instrumento directamente negociable y replicante.


\section{Construcción del conjunto de datos sintético}
\label{sec:synthetic_dataset}

No existe una base observable de opciones asiáticas líquidas sobre H100 que
proporcione simultáneamente precios y Greeks. La serie real se utiliza para
caracterizar la dinámica física y definir regiones plausibles; las etiquetas de
precio y Delta se generan mediante RQMC bajo las leyes de valoración del
experimento. Por tanto, los targets son datos sintéticos de modelo.

Con $K=1$, se define $s_t=P_t/K$ y
\begin{equation}
    c_t=\frac{C_t}{K}
    =\frac{1}{N}\sum_{i=1}^{m(t)}\frac{P_{t_i}}{K}.
\end{equation}
Si $a_t$ es la media de los fixings ya realizados, entonces
$c_t=(m(t)/N)a_t$. El vector final de inputs es
\begin{equation}
    \mathbf x=\left(s_t,c_t,\tau,r,\kappa_Q,\theta_Q,\sigma_Q\right)
    \in\mathbb R^7.
    \label{eq:final_ml_inputs}
\end{equation}
Los parámetros de $\mathbb Q$ se incluyen para que un único surrogate aprenda
una familia de funciones de valoración.

\begin{table}[htbp]
    \centering
    \caption{Dominio del experimento sintético Log-OU.}
    \label{tab:synthetic_domain}
    \begin{tabular}{lll}
        \toprule
        Variable & Rango & Papel \\
        \midrule
        $P_0/K$ & $[0{,}85,1{,}175]$ & condición inicial física \\
        $\tau$ & $[0{,}05,1{,}00]$ & vencimiento restante \\
        $r$ & $[0{,}00,0{,}08]$ & descuento \\
        $\kappa_Q$ & $[25,87]$ & reversión bajo $\mathbb Q$ \\
        $\theta_Q$ & $[-0{,}1641,0{,}1610]$ & nivel Log-OU bajo $\mathbb Q$ \\
        $\sigma_Q$ & $[0{,}55,0{,}71]$ & volatilidad bajo $\mathbb Q$ \\
        $P_t/K$ & $[0{,}70,1{,}30]$ & filtro del estado realizado \\
        $a_t$ & $[0{,}70,1{,}30]$ & filtro de media pasada \\
        \bottomrule
    \end{tabular}
\end{table}

Un diseño Sobol aleatorizado cubre inicialmente las seis variables exógenas
\[
    (P_0/K,\tau,r,\kappa_Q,\theta_Q,\sigma_Q).
\]
Una vez fijado
$t=T-\tau$, se genera bajo $\mathbb P$ una única historia Log-OU hasta $t$ y de
ella se obtienen conjuntamente $P_t/K$ y $c_t$. Ambos estados no se muestrean de
forma independiente. Los candidatos se aceptan cuando el spot realizado y, si
existen fixings previos, su media se encuentran dentro del dominio de la
Tabla~\ref{tab:synthetic_domain}. Después, condicionado ese estado, los fixings
futuros y las etiquetas $(V,\Delta)$ se generan bajo $\mathbb Q$. El esquema es,
por tanto,
\begin{equation}
    \begin{aligned}
        \text{diseño paramétrico}
        &\rightarrow \text{historia bajo }\mathbb P
        \rightarrow (P_t,C_t) \\
        &\rightarrow \text{valoración bajo }\mathbb Q .
    \end{aligned}
    \label{eq:dataset_measure_pipeline}
\end{equation}

Los tamaños congelados del diseño experimental son $65\,536$ estados de entrenamiento,
$8\,192$ de validación, $8\,192$ de test y $1\,024$ de referencia. Para estudiar
eficiencia muestral, el training se organiza en subconjuntos anidados
\begin{equation}
    \mathcal D_{1024}\subset\mathcal D_{4096}\subset
    \mathcal D_{16384}\subset\mathcal D_{65536}.
    \label{eq:nested_training_sets}
\end{equation}
Los cuatro conjuntos principales se generan independientemente desde el origen,
con semillas diferentes para el estado físico y para la valoración RQMC.

El número de fixings futuros y la distancia al siguiente fixing se conservan
como metadatos, pero no se incorporan al vector de siete inputs. Esta decisión se
mantuvo congelada antes del test. Con el calendario fijo, $\tau$ determina el número de fixings pendientes y sus
distancias temporales; con la convención posterior al fixing,
$m(t)=\lfloor12t\rfloor$ para $0\leq t<1$, salvo tolerancias numéricas.
Incluir esos metadatos podría facilitar la aproximación de los cambios de régimen,
aunque no añadiría información al estado. En cada fixing, $C_t$ incluye ya el precio observado.


\section{MLP supervisado únicamente por precios}
\label{sec:mlp}

El benchmark neuronal es una MLP que aprende el precio normalizado
$v=V/K$ únicamente a partir de las etiquetas de valor. Las redes feedforward son
aproximadores flexibles de funciones multivariantes \parencite{hornik1989}; su
calidad se determina aquí de forma empírica sobre muestras no utilizadas para el
ajuste.

Inputs y precio se estandarizan con medias y desviaciones calculadas
exclusivamente sobre cada subconjunto de entrenamiento,
\begin{equation}
    \widetilde x_j=\frac{x_j-\mu_{x_j}}{\sigma_{x_j}},\qquad
    \widetilde v=\frac{v-\mu_v}{\sigma_v}.
    \label{eq:input_standardization}
\end{equation}
La arquitectura compartida por MLP y DML es residual, con anchura 128, activación
SiLU y tres bloques residuales de dos transformaciones lineales, siguiendo la
idea general de aprendizaje residual \parencite{he_etal_2016}. Con siete inputs
contiene $100\,225$ parámetros entrenables. La salida es lineal y aproxima el
precio estandarizado.

El MLP minimiza únicamente
\begin{equation}
    \mathcal L_{MLP}
    =\frac{1}{n}\sum_{i=1}^{n}
    \left(f_{\omega}(\widetilde{\mathbf x}_i)-\widetilde v_i\right)^2,
    \label{eq:mlp_loss}
\end{equation}
por lo que la Delta no interviene en el entrenamiento. La optimización sigue la Tabla~\ref{tab:shared_ml_configuration}.
Se ejecutan 3000 épocas sin \textit{early stopping} y se conserva el checkpoint
con menor MSE de precio estandarizado en validación.

Aunque no recibe supervisión diferencial, la red es diferenciable. Mediante
autodiferenciación \parencite{baydin2018}, su Delta implícita es
\begin{equation}
    \widehat\Delta^{MLP}
    =\frac{\sigma_v}{\sigma_s}
    \frac{\partial f_\omega}{\partial\widetilde s},
    \label{eq:mlp_autograd_delta}
\end{equation}
donde $s=P_t/K$. Así puede compararse la calidad de una derivada aprendida de
forma implícita con la obtenida cuando esa derivada forma parte explícita de la
supervisión.


\section{Differential Machine Learning}
\label{sec:dml}

DML amplía la regresión por valores incorporando información sobre derivadas de
la función objetivo, siguiendo el enfoque de \textcite{huge_savine_2020}. Para
aislar el efecto de esta información, DML utiliza exactamente la misma
arquitectura, inicialización, optimizador y procedimiento de checkpoint que el
MLP. La única diferencia sistemática es la función de pérdida.

La Delta debe transformarse de forma coherente con la estandarización. Si
$\widetilde s=(s-\mu_s)/\sigma_s$ y
$\widetilde v=(v-\mu_v)/\sigma_v$, la etiqueta diferencial en coordenadas
estandarizadas es
\begin{equation}
    \widetilde\Delta
    =\frac{\partial\widetilde v}{\partial\widetilde s}
    =\Delta\frac{\sigma_s}{\sigma_v}.
    \label{eq:dml_scaled_delta}
\end{equation}
Para evitar que su escala domine mecánicamente la función objetivo se aplica una
normalización RMS calculada únicamente sobre training,
\begin{equation}
    w_\Delta=\left[\frac{1}{n}\sum_{i=1}^{n}
    \widetilde\Delta_i^2\right]^{-1/2}.
    \label{eq:dml_derivative_weight}
\end{equation}

La pérdida combina precio y derivada,
\begin{equation}
    \mathcal L_{DML}=\alpha\mathcal L_V+\beta\mathcal L_\Delta,
    \qquad
    \alpha=\frac{1}{1+\lambda_\Delta},\quad
    \beta=\frac{\lambda_\Delta}{1+\lambda_\Delta},
    \label{eq:dml_total_loss}
\end{equation}
donde $\mathcal L_V$ es el MSE del precio estandarizado y
$\mathcal L_\Delta$ el MSE de la derivada estandarizada y ponderada por
$w_\Delta$. La derivada de la salida neuronal respecto al spot se obtiene por
autodiferenciación dentro del propio grafo de entrenamiento, por lo que el
backpropagation de $\mathcal L_\Delta$ requiere derivadas de orden superior
respecto a los parámetros de la red.

En el protocolo final, el peso seleccionado antes del test es $\lambda_\Delta=0{,}25$, equivalente a $\alpha=0{,}8$ y $\beta=0{,}2$.
DML no recibe supervisión sobre otras Greeks ni modifica la arquitectura para
acomodar la información adicional. Por tanto, la comparación evalúa de forma
directa si añadir la Delta a cada observación mejora la aproximación del precio
y, especialmente, de su gradiente.


\section{Selección de hiperparámetros y protocolo de congelación}
\label{sec:model_selection}

Para evitar que la comparación MLP--DML quede condicionada por una búsqueda
arquitectónica específica del problema de \textit{compute}, la arquitectura y el
optimizador se heredan del benchmark GBM previamente confirmado y congelado. En
ese benchmark, la selección se realizó sobre un MLP entrenado sólo con precios
mediante una regla de un error estándar sobre el MSE de precio medio
multi-seed; entre los candidatos admisibles se priorizaron después menor número
de parámetros y menor tiempo de entrenamiento. La Tabla~\ref{tab:shared_ml_configuration} recoge la configuración transferida al
problema Log-OU. No se reoptimizan arquitectura ni optimizador tras
introducir los dos parámetros adicionales del estado Log-OU; con los siete inputs
finales la red contiene $100\,225$ parámetros entrenables.

\begin{table}[htbp]
    \centering
    \caption{Configuración neuronal compartida por MLP y DML.}
    \label{tab:shared_ml_configuration}
    \begin{tabular}{ll}
        \toprule
        Hiperparámetro & Configuración final \\
        \midrule
        Backbone & Residual \\
        Anchura / bloques & 128 / 3 \\
        Activación & SiLU \\
        Optimizador / scheduler & AdamW / coseno \\
        Learning rate inicial & $1{,}9616\times10^{-3}$ \\
        Weight decay & $6{,}6341\times10^{-7}$ \\
        Batch size & 1024 \\
        Máximo de épocas & 3000 \\
        Checkpoint & MSE de precio estandarizado \\
        \bottomrule
    \end{tabular}
\end{table}

Para el problema de \textit{compute} sólo se vuelve a seleccionar el peso
diferencial. Se evalúa
$\lambda_\Delta\in\{0{,}25,0{,}5,1,2,4\}$ con $8192$ observaciones y las semillas
$\{1701,3701,5701\}$, con un máximo de 1500 épocas. Tras una permutación
determinista, el conjunto de validación destinado a esta etapa se divide en dos
mitades disjuntas: la primera selecciona el epoch exclusivamente mediante MSE de
precio estandarizado; la segunda compara los valores de $\lambda_\Delta$ y nunca
interviene en la selección del checkpoint. El criterio de comparación es
\begin{equation}
    \mathcal S_{\lambda}
    =\frac{1}{2}\,\mathrm{MSE}_{V,\mathrm{std}}
    +\frac{1}{2}\,\mathrm{MSE}_{\Delta,\mathrm{std,RMS}},
    \label{eq:lambda_selection_score}
\end{equation}
y se minimiza su media entre semillas. El valor seleccionado es
$\lambda_\Delta=0{,}25$. El valor $4{,}0$ pertenece al benchmark previo y no se
utiliza como peso diferencial del experimento final de \textit{compute}.

La secuencia experimental puede resumirse como
\begin{equation}
    \begin{aligned}
        \text{benchmark GBM congelado}
        &\rightarrow \text{transferencia al estado Log-OU} \\
        &\rightarrow \text{selección de }\lambda_\Delta
        \rightarrow \textit{freeze}
        \rightarrow \text{evaluación final}.
    \end{aligned}
    \label{eq:model_selection_pipeline}
\end{equation}
Antes de evaluar test y referencia quedan fijados los siete inputs, arquitectura,
peso diferencial, normalización, batch size, criterio de checkpoint, tamaños y
semillas. Los diagnósticos posteriores sólo se utilizan para interpretar errores
y limitaciones; no permiten reentrenar ni modificar el protocolo.

La evaluación final emplea cuatro tamaños de entrenamiento: $1\,024$, $4\,096$,
$16\,384$ y $65\,536$ observaciones. Para cada tamaño se utilizan cinco semillas:
$1701$, $2701$, $3701$, $4701$ y $5701$. Para cada pareja tamaño--semilla se
entrenan un MLP y un DML con la misma arquitectura y protocolo, generando 40 ejecuciones. Las muestras
de entrenamiento son anidadas, de modo que la curva de aprendizaje compara
cantidades crecientes de información manteniendo la estructura experimental.
Las cinco semillas varían el entrenamiento sobre los mismos estados de evaluación;
su dispersión no mide incertidumbre por regeneración del dataset. Los veinte pares
tamaño--semilla tampoco constituyen veinte réplicas independientes del diseño de datos.

\section{Validación numérica, consistencia y control de leakage}
\label{sec:validation}

La validación se estructura alrededor de cuatro riesgos: coherencia de los
estados, corrección de las etiquetas, consistencia computacional y aislamiento
del test. Primero se comprueba que cada escenario satisface la identidad
contractual $n_{past}+n_{future}=N$, que $C_t$ coincide con la contribución de los
fixings ya observados, que los parámetros y estados están dentro del dominio y
que no existen valores no finitos.

Train, validación, test y referencia se generan con identificadores y semillas
separados para la historia física y para el RQMC. Se comprueba la ausencia de
solapamientos entre escenarios y aleatorizaciones. El scaler de inputs y precio,
así como el factor RMS diferencial, se ajustan exclusivamente sobre el training
de cada ejecución; las demás muestras se transforman con esos parámetros ya
fijados.

Los controles conservados abarcan 12 estados para diferencias finitas centrales
con \textit{common random numbers} y otros 12 para la comparación CPU/Torch.
La discrepancia pathwise/FD media es $1{,}8811\times10^{-7}$ y la máxima
$1{,}3414\times10^{-6}$; los máximos CPU/Torch son $5{,}5511\times10^{-17}$
en precio y $6{,}9389\times10^{-18}$ en Delta. Son controles puntuales satisfactorios,
no una cota uniforme del error de las etiquetas. El dataset final y el Oracle
dinámico usan float32; la coincidencia a precisión máquina del control conservado
no certifica todo ese recorrido. La prima inicial dinámica agrega ocho
aleatorizaciones de 8192 puntos en float64 y el Oracle dinámico usa 4096 puntos
por estado y una aleatorización.

Los metadatos finales registran PyTorch 2.12.1+cu126 y CUDA 12.6. El material
reproducible distingue estas versiones verificadas de las dependencias compatibles
y conserva la procedencia de los scripts recuperados para diagnósticos intermedios.

Finalmente, el protocolo queda congelado antes del test y cada ejecución se
vincula a su configuración. MLP y DML se comparan de manera emparejada, con la
misma arquitectura, mismos estados, misma inicialización y mismo orden de
minibatches; la diferencia deliberada es la supervisión de Delta. Los listados
completos de semillas, hashes y comprobaciones operativas se reservan para el
material reproducible, ya que no alteran la lógica estadística del experimento.


\section{Diseños de cobertura y evaluación económica}
\label{sec:hedging_design}

La precisión de una Greek no garantiza por sí sola una mejora económica. Por
ello se plantean dos ejercicios complementarios: una cobertura local, que aísla
la calidad de la sensibilidad, y un stress test dinámico con un forward de
cómputo modelizado, que incorpora la conversión de Delta en una decisión de
cobertura.

\subsection*{Cobertura local}

Para un estado $(P_t,C_t,\ldots)$ se aplica un pequeño shock
$\Delta P=P_t\varepsilon$ manteniendo fija la parte histórica del contrato. El
cambio real de valor se estima con RQMC y se compara con la aproximación de
primer orden $\widehat\Delta_t\Delta P$. El residuo cubierto es
\begin{equation}
    R^{m}=\Delta V-\widehat\Delta_t^{m}\Delta P,
    \qquad m\in\{MLP,DML,oracle\},
    \label{eq:local_hedge_residual}
\end{equation}
y el benchmark sin cobertura es $R^0=\Delta V$. El shock principal es del 1\%,
con robustez al 0{,}5\% y 2\%, en ambas direcciones. Se evalúan RMSE, MAE y
comparaciones emparejadas. Este ejercicio mide directamente la capacidad de la
Delta para neutralizar pequeñas variaciones del valor, sin convertirla todavía
en una estrategia dinámica. Se emplean 6140 estados procedentes de cuatro bloques
de 128 trayectorias, con dos direcciones por shock: 12\,280 comparaciones dependientes.
Los estados se toman en el punto medio de los primeros once intervalos y en un
punto interior del último, a unos $15{,}21$ y $8{,}52$ días de la frontera más próxima,
respectivamente. La prueba dinámica evalúa, en cambio, estados posteriores a cada fixing.

\subsection*{Cobertura dinámica con forward modelizado}

El spot de capacidad no se considera directamente negociable en la forma
necesaria para una réplica clásica. Se introduce, por tanto, un forward
conceptual que liquida en el siguiente fixing. Bajo el escenario $Q_0$, si
$h_j=t_{j+1}-t_j$, la transición Log-OU implica
\begin{equation}
    F_j=\mathbb E^{Q_0}[P_{t_{j+1}}\mid\mathcal F_{t_j}]
    =\exp\left[\theta_Q+a_j(X_{t_j}-\theta_Q)+\frac{1}{2}\nu_j^2\right],
    \qquad a_j=e^{-\kappa_Qh_j},
    \label{eq:logou_forward}
\end{equation}
donde $\nu_j^2=\sigma_Q^2(1-e^{-2\kappa_Qh_j})/(2\kappa_Q)$. Su sensibilidad
al spot es
\begin{equation}
    J_j=\frac{\partial F_j}{\partial P_{t_j}}
    =e^{-\kappa_Qh_j}\frac{F_j}{P_{t_j}},
    \label{eq:forward_spot_jacobian}
\end{equation}
y la política efectivamente ejecutada para una opción vendida fue
\begin{equation}
    q_j=\frac{\widehat\Delta_{t_j}}{J_j}.
    \label{eq:forward_hedge_ratio}
\end{equation}
Esta política histórica omite el descuento de la sensibilidad del contrato.
Para un forward con precio de entrega fijado $L$, el valor actual y la posición
que neutraliza exactamente la Delta del valor de la opción son
\begin{equation}
    G_t=e^{-r(u-t)}[F(t,u)-L],\qquad
    q_t^{\mathrm{PV}}=e^{r(u-t)}\frac{\widehat\Delta_t}{J_t}.
    \label{eq:forward_pv_correction}
\end{equation}
Se distingue el precio de entrega $F$ del valor del contrato $G$: coste inicial
nulo no permite omitir el descuento. Las tablas principales conservan la política
congelada $q_j=\widehat\Delta/J_j$; la Sección~\ref{sec:dynamic_hedging}
cuantifica por separado la corrección posterior, sin reentrenamiento.
La división por $J_j$ es esencial: con reversión rápida, la sensibilidad del
forward al spot puede ser reducida y amplificar el impacto de un error de Delta
sobre el nocional de cobertura.

La estrategia se rebalancea después de cada fixing. Las trayectorias económicas
se generan bajo $\mathbb P$, mientras que precio, Deltas y forward se calculan
bajo $Q_0$. El forward contratado en $t_j$ produce en $t_{j+1}$
$q_j(P_{t_{j+1}}-F_j)$. Desde la perspectiva del vendedor de la opción, el error
terminal de una estrategia $m$ es
\begin{equation}
    E_T^{m}=V_0e^{rT}
    +\sum_{j=0}^{N-1}q_j^{m}(P_{t_{j+1}}-F_j)e^{r(T-t_{j+1})}
    -H_T,
    \label{eq:dynamic_terminal_hedge_error}
\end{equation}
y sin cobertura $E_T^0=V_0e^{rT}-H_T$.

El stress test central utiliza $P_0/K=1$, $r=3\%$, parámetros de $Q_0$ y cuatro
bloques independientes de 512 trayectorias físicas. Sólo se conservan trayectorias
cuyos estados de rebalanceo permanecen dentro del dominio común de los modelos:
2038 de 2048, el $99{,}51\%$. El filtro utiliza la permanencia en todos los
rebalanceos y es común a las estrategias; no selecciona por sus errores.
Las métricas son condicionales a esa permanencia, no el riesgo incondicional bajo
$\mathbb P$. Las decisiones de cobertura usan información contemporánea.
No se aplica clipping a Deltas ni a nocionales. MLP y DML utilizan el promedio de
los cinco modelos finales entrenados con $n=65\,536$, mientras que el oracle
recalcula la Delta por RQMC.

La evaluación dinámica utiliza RMSE y MAE del error terminal, percentiles de
$|E_T|$, concentración del error cuadrático en las colas y comparaciones
emparejadas frente a MLP y no hedge. Además, se diagnostica el error por número
de fixings pendientes y proximidad a la siguiente fecha contractual, ya que esas
variables discretas no forman parte explícita del estado neuronal.

La prueba local evalúa $dV\simeq\Delta\,dP$ en estados interiores y la dinámica
evalúa la política forward en fronteras, ausentes exactamente del entrenamiento
sintético. Cambian simultáneamente calendario, instrumento y acumulación de errores;
su comparación no aísla causalmente cada mecanismo.
